\documentclass[11pt,letterpaper]{article}
\usepackage[margin=0.75in]{geometry}
\usepackage{fancyhdr}
\usepackage{tcolorbox}
\usepackage{amsmath}
\usepackage{amssymb}
\usepackage{enumitem}
\usepackage{graphicx}
\usepackage[hidelinks]{hyperref}
\usepackage{tikz}

\pagestyle{fancy}
\fancyhf{}
\fancyhead[L]{\textbf{Lesson 2: Schematics \& Troubleshooting}}
\fancyhead[R]{\textbf{Page \thepage}}
\renewcommand{\headrulewidth}{2pt}
\setlength{\headheight}{14pt}

\tcbuselibrary{skins,breakable}

% Custom boxes
\newtcolorbox{conceptbox}{
  colback=blue!5,
  colframe=blue!75!black,
  fonttitle=\bfseries,
  breakable,
  enhanced jigsaw
}

\newtcolorbox{warningbox}{
  colback=red!5,
  colframe=red!75!black,
  fonttitle=\bfseries,
  title=WARNING,
  breakable,
  enhanced jigsaw
}

\newtcolorbox{examplebox}{
  colback=green!5,
  colframe=green!75!black,
  fonttitle=\bfseries,
  breakable,
  enhanced jigsaw
}

\newtcolorbox{keyinsight}{
  colback=yellow!10,
  colframe=orange!80!black,
  fonttitle=\bfseries,
  title=KEY INSIGHT,
  breakable,
  enhanced jigsaw
}

\setlength{\parindent}{0pt}
\setlength{\parskip}{0.5em}

\begin{document}

\begin{center}
{\Huge \textbf{LESSON 2}}\\[0.5em]
{\LARGE \textbf{Reading Schematics, Datasheets, and}}\\
{\LARGE \textbf{Practical Troubleshooting}}\\[0.5em]
{\large NLC STEM Club UAV Project}\\
{\large Electrical Systems Education Series}\\[0.5em]
\rule{\textwidth}{2pt}
\end{center}

\begin{conceptbox}
\textbf{PREREQUISITES:}
\begin{itemize}[nosep]
\item Completed Lesson 1: Circuit Fundamentals Teaching Guide
\item Understanding of Ohm's Law, power calculations, and wire gauge selection
\item Basic familiarity with your project components
\end{itemize}
\end{conceptbox}

\begin{conceptbox}
\textbf{LESSON OBJECTIVES:}

By the end of this lesson, you will be able to:
\begin{enumerate}[nosep]
\item Read and interpret electrical schematics and wiring diagrams
\item Extract critical information from component datasheets
\item Use a multimeter to measure voltage, current, and continuity
\item Understand PWM signals and how they control servos and ESCs
\item Troubleshoot common electrical problems systematically
\item Recognize and prevent grounding issues and electrical noise
\item Apply these skills to debug your UAV electrical system
\end{enumerate}
\end{conceptbox}

\newpage
\section{Reading Electrical Schematics}

\subsection{What is a Schematic?}

A schematic is a symbolic representation of an electrical circuit that shows:
\begin{itemize}
\item How components are electrically connected
\item The flow of current and signals
\item Component values and part numbers
\end{itemize}

\textbf{KEY DIFFERENCE: Schematic vs Wiring Diagram}

\begin{tabular}{|p{0.45\textwidth}|p{0.45\textwidth}|}
\hline
\textbf{SCHEMATIC} & \textbf{WIRING DIAGRAM} \\
\hline
Shows ELECTRICAL connections (how electricity flows) & Shows PHYSICAL connections (which wire goes where) \\
\hline
Uses standardized symbols & Shows actual connectors and wire colors \\
\hline
Does NOT show physical layout & Shows spatial relationships between components \\
\hline
Used by engineers to design and analyze circuits & Used during assembly and troubleshooting \\
\hline
\end{tabular}

\subsection{Common Schematic Symbols}

\textbf{POWER SOURCES:}

\begin{itemize}
\item \textbf{Battery:} Two parallel lines (longer line is positive)
\item \textbf{DC Power Supply:} Box with V+ and GND labels
\end{itemize}

\textbf{CONNECTIONS:}

\begin{itemize}
\item \textbf{Wire connection:} Dot indicates electrical connection
\item \textbf{Wires crossing:} No dot = no connection, just crossing
\item \textbf{Ground symbols:} Earth ground, chassis ground, signal/digital ground
\end{itemize}

\textbf{PASSIVE COMPONENTS:}

\begin{itemize}
\item \textbf{Resistor:} Zigzag line or rectangle
\item \textbf{Capacitor:} Two parallel lines (electrolytic has + polarity)
\item \textbf{Inductor:} Coil symbol
\end{itemize}

\textbf{ACTIVE COMPONENTS:}

\begin{itemize}
\item \textbf{Diode:} Triangle with line (current flows in arrow direction)
\item \textbf{LED:} Diode with small arrows pointing out
\item \textbf{Transistor (NPN):} Three terminals: Base, Collector, Emitter
\item \textbf{Voltage Regulator:} Box with IN, OUT, and GND pins
\end{itemize}

\textbf{MOTORS \& MECHANICAL:}

\begin{itemize}
\item \textbf{Motor (DC):} Circle with M inside
\item \textbf{Motor (Servo):} Box with mechanical coupling
\end{itemize}

\textbf{SWITCHES \& SENSORS:}

\begin{itemize}
\item \textbf{Switch:} SPST (Single Pole Single Throw)
\item \textbf{Button:} Normally open pushbutton
\end{itemize}

\subsection{Reading a Schematic: Step-by-Step}

\textbf{STEP 1: Identify the Power Sources}
\begin{itemize}
\item Look for battery symbols or ``V+'' labels
\item Note the voltage level (e.g., 14.8V, 5V, 3.3V)
\item Locate all ground symbols
\end{itemize}

\textbf{STEP 2: Follow the Power Rails}
\begin{itemize}
\item Trace from positive terminal through components to ground
\item Identify parallel branches (components sharing same voltage)
\item Identify series components (current flows through both)
\end{itemize}

\textbf{STEP 3: Identify Signal Paths}
\begin{itemize}
\item Find input signals (from sensors, receivers)
\item Trace signal wires to processing units (Arduino, flight controller)
\item Find output signals (to ESC, servos, LEDs)
\end{itemize}

\textbf{STEP 4: Note Component Values}
\begin{itemize}
\item Resistor values ($\Omega$)
\item Capacitor values ($\mu$F, pF)
\item Component part numbers or ratings
\end{itemize}

\begin{examplebox}
\textbf{EXAMPLE - Simple LED Circuit:}

Current flows: 5V $\rightarrow$ through R1 (220$\Omega$) $\rightarrow$ through LED $\rightarrow$ to GND

\textbf{HOW TO READ THIS:}
\begin{enumerate}[nosep]
\item Power source: 5V
\item Current flows from 5V through resistor through LED to ground
\item Resistor limits current to safe level for LED
\item Voltage drop: 5V - 2V (LED) = 3V across resistor
\item Current: $I = V/R = 3V / 220\Omega = 13.6$ mA
\end{enumerate}
\end{examplebox}

\subsection{UAV System Schematic Example}

Let's analyze a simplified version of YOUR dual-BEC system:

\begin{center}
\texttt{[4S LiPo Battery: 14.8V]}\\
$\downarrow$\\
\texttt{XT60 Y-Harness}\\
$\swarrow$ \hspace{3cm} $\searrow$\\
\texttt{ESC} \hspace{3cm} \texttt{External UBEC}\\
\texttt{5V @ 3A} \hspace{2.5cm} \texttt{5V @ 6A}\\
\texttt{RAIL 1} \hspace{2.5cm} \texttt{RAIL 2}\\
$\downarrow$ \hspace{4cm} $\downarrow$\\
Flight Critical \hspace{1.5cm} Servos\\
Components
\end{center}

\begin{keyinsight}
This schematic answers ``WHY two BECs?''
\begin{itemize}
\item Total servo current could exceed ESC BEC rating
\item Separating rails prevents servo noise from affecting flight controller
\item If servos brown out Rail 2, Rail 1 stays stable
\end{itemize}
\end{keyinsight}

\newpage
\section{Reading Datasheets}

\subsection{What is a Datasheet?}

A datasheet is a technical document provided by the manufacturer that contains:
\begin{itemize}
\item Electrical specifications (voltage, current, power)
\item Mechanical dimensions and pinouts
\item Performance characteristics (speed, accuracy, range)
\item Operating conditions (temperature, humidity)
\item Application guidelines and example circuits
\end{itemize}

\textbf{DATASHEETS ARE ESSENTIAL FOR:}
\begin{itemize}
\item Verifying compatibility with your system
\item Determining proper operating conditions
\item Connecting components correctly (pinouts)
\item Troubleshooting when things don't work
\end{itemize}

\subsection{Key Sections to Look For}

\textbf{SECTION 1: ABSOLUTE MAXIMUM RATINGS}

\begin{warningbox}
CRITICAL - exceeding these can destroy the component!

Example (from a typical voltage regulator datasheet):
\begin{itemize}[nosep]
\item Input Voltage: 4.5V to 40V
\item Output Current: 3A continuous, 5A peak
\item Operating Temperature: -40$^\circ$C to +125$^\circ$C
\item Storage Temperature: -65$^\circ$C to +150$^\circ$C
\end{itemize}

If you apply 50V to this regulator = instant damage!
\end{warningbox}

\textbf{SECTION 2: RECOMMENDED OPERATING CONDITIONS}

These are SAFE ranges for normal operation

Example:
\begin{itemize}
\item Input Voltage: 7V to 30V (within absolute max)
\item Output Current: 0 to 3A (with heatsink if $>$ 1.5A)
\item Operating Temperature: 0$^\circ$C to 85$^\circ$C
\end{itemize}

Design your system to stay within these ranges!

\textbf{SECTION 3: ELECTRICAL CHARACTERISTICS}

Detailed performance specifications

\begin{center}
\begin{tabular}{|l|l|c|c|c|c|}
\hline
\textbf{Parameter} & \textbf{Condition} & \textbf{Min} & \textbf{Typ} & \textbf{Max} & \textbf{Unit} \\
\hline
Output Voltage & Vout = 5V & 4.9 & 5.0 & 5.1 & V \\
\hline
Line Regulation & Vin = 7V-30V & - & 10 & 50 & mV \\
\hline
Load Regulation & Iout = 0-3A & - & 20 & 100 & mV \\
\hline
Quiescent Current & No load & - & 10 & 15 & mA \\
\hline
Efficiency & Vin=12V, Iout=1A & - & 85 & - & \% \\
\hline
\end{tabular}
\end{center}

\textbf{HOW TO READ THIS:}
\begin{itemize}
\item ``Typ'' = typical value (what you'll usually measure)
\item ``Min/Max'' = guaranteed range (your unit will fall within this)
\item Output voltage will be $5.0V \pm 0.1V$ (4.9V to 5.1V)
\end{itemize}

\textbf{SECTION 4: PIN CONFIGURATION (PINOUT)}

Shows which pin does what

\begin{warningbox}
Connecting wrong pins = component destruction!
\end{warningbox}

\textbf{SECTION 5: APPLICATION CIRCUIT}

Example schematic showing proper use. Often includes:
\begin{itemize}
\item Input/output capacitors (for stability)
\item Protection diodes
\item LED indicators
\item Recommended wire gauges
\end{itemize}

\textbf{SECTION 6: PACKAGE INFORMATION}

Physical dimensions, hole spacing, mounting

\subsection{Real-World Datasheet Example: BME280 Sensor}

Let's extract key information for YOUR air quality sensor project:

\textbf{FROM THE BME280 DATASHEET:}

\begin{tcolorbox}[title=ABSOLUTE MAXIMUM RATINGS]
\begin{itemize}[nosep]
\item Supply Voltage (Vdd): -0.3V to 4.25V
\item ESD Tolerance: 2000V
\end{itemize}
$\rightarrow$ Connect to Arduino 3.3V output (NOT 5V!) to stay safe OR use a level shifter module
\end{tcolorbox}

\begin{tcolorbox}[title=OPERATING CONDITIONS]
\begin{itemize}[nosep]
\item Supply Voltage: 1.71V to 3.6V
\item Supply Current: 3.6 $\mu$A @ 1 Hz (typical), 714 $\mu$A @ 1 Hz (max)
\end{itemize}
$\rightarrow$ Very low power! Won't strain your Arduino power budget
\end{tcolorbox}

\begin{tcolorbox}[title=MEASUREMENT RANGES \& ACCURACY]
\begin{itemize}[nosep]
\item Temperature: -40$^\circ$C to +85$^\circ$C ($\pm$1$^\circ$C accuracy)
\item Humidity: 0\% to 100\% RH ($\pm$3\% accuracy)
\item Pressure: 300 to 1100 hPa ($\pm$1 hPa accuracy)
\end{itemize}
$\rightarrow$ Perfect for your environmental monitoring mission!\\
$\rightarrow$ Can detect altitude changes of $\sim$8 meters
\end{tcolorbox}

\begin{tcolorbox}[title=I2C INTERFACE]
\begin{itemize}[nosep]
\item I2C Address: 0x76 or 0x77 (selectable)
\item Clock Speed: Up to 3.4 MHz
\item Logic Levels: VddIO (0.3 $\times$ Vdd to Vdd)
\end{itemize}
$\rightarrow$ Use Wire.begin() in Arduino code\\
$\rightarrow$ Check address with I2C scanner if not responding
\end{tcolorbox}

\begin{tcolorbox}[title=PINOUT (8-pin package)]
\begin{itemize}[nosep]
\item Pin 1: GND - Ground
\item Pin 2: CSB - Chip Select (I2C: connect to Vdd)
\item Pin 3: SDI - I2C Data (SDA)
\item Pin 4: SCK - I2C Clock (SCL)
\item Pin 5: SDO - Address select (GND=0x76, Vdd=0x77)
\item Pin 6: Vdd - Power supply 1.71V-3.6V
\item Pin 7: VddIO - Logic level reference (can = Vdd)
\item Pin 8: GND - Ground
\end{itemize}
\end{tcolorbox}

\textbf{WIRING TO ARDUINO:}

\begin{center}
\begin{tabular}{ll}
\textbf{BME280 Breakout Board} & \textbf{Arduino Nano} \\
\hline
VCC/Vin & 3.3V (NOT 5V!) \\
GND & GND \\
SDA & A4 (I2C Data) \\
SCL & A5 (I2C Clock) \\
\end{tabular}
\end{center}

\subsection{Practice: Extract Info from Your Component Datasheets}

For EACH component in your system, find and record:

\textbf{ESC DATASHEET:}
\begin{enumerate}
\item Continuous current rating: \underline{\hspace{2cm}} A
\item Input voltage range: \underline{\hspace{2cm}} V to \underline{\hspace{2cm}} V
\item BEC output voltage: \underline{\hspace{2cm}} V
\item BEC output current: \underline{\hspace{2cm}} A
\item PWM signal frequency: \underline{\hspace{2cm}} Hz
\item Wiring: Red = \underline{\hspace{2cm}}, Black = \underline{\hspace{2cm}}, White/Yellow = \underline{\hspace{2cm}}
\end{enumerate}

\textbf{UBEC DATASHEET:}
\begin{enumerate}
\item Input voltage range: \underline{\hspace{2cm}} V to \underline{\hspace{2cm}} V
\item Output voltage: \underline{\hspace{2cm}} V (adjustable? Y/N)
\item Continuous output current: \underline{\hspace{2cm}} A
\item Peak output current: \underline{\hspace{2cm}} A
\item Efficiency at your voltage: \underline{\hspace{2cm}} \%
\end{enumerate}

\textbf{FLIGHT CONTROLLER (APM 2.8) DATASHEET:}
\begin{enumerate}
\item Input voltage: \underline{\hspace{2cm}} V (connect to which BEC?)
\item PWM output voltage: \underline{\hspace{2cm}} V
\item Number of output channels: \underline{\hspace{2cm}}
\item GPS connector pinout: TX = pin \underline{\hspace{1cm}}, RX = pin \underline{\hspace{1cm}}
\item Telemetry port: \underline{\hspace{3cm}}
\end{enumerate}

\newpage
\section{Using a Multimeter}

\subsection{Multimeter Basics}

A multimeter measures three fundamental electrical properties:
\begin{itemize}
\item \textbf{VOLTAGE (V or VDC/VAC)} - Potential difference
\item \textbf{CURRENT (A or mA)} - Flow of electrons
\item \textbf{RESISTANCE ($\Omega$)} - Opposition to current
\end{itemize}

PLUS useful functions:
\begin{itemize}
\item \textbf{CONTINUITY TEST} (beep when connected)
\item \textbf{DIODE TEST} (measures forward voltage drop)
\end{itemize}

\textbf{DIAL POSITIONS (typical multimeter):}
\begin{itemize}
\item \textbf{OFF} - Power off
\item \textbf{V$\sim$ (ACV)} - AC voltage (wall outlets: 120V)
\item \textbf{V- (DCV)} - DC voltage $\leftarrow$ YOU'LL USE THIS MOST
\item \textbf{$\Omega$} - Resistance
\item \textbf{Continuity} - Continuity test (beeper)
\item \textbf{Diode} - Diode test
\item \textbf{A} - Current (often separate jack!)
\end{itemize}

\textbf{PROBE PLACEMENT:}
\begin{itemize}
\item RED probe: Positive terminal, higher voltage point
\item BLACK probe: Negative terminal, ground, lower voltage point
\end{itemize}

\subsection{Measuring Voltage}

Voltage is measured IN PARALLEL (across a component)

\textbf{PROCEDURE:}
\begin{enumerate}
\item Set dial to ``V-'' (DC voltage)
\item Select range: 20V for most UAV measurements
\item Touch RED probe to positive point
\item Touch BLACK probe to ground/negative point
\item Read the display
\end{enumerate}

\begin{examplebox}
\textbf{EXAMPLE 1: Checking Battery Voltage}

Battery terminals: (+) and (-)\\
Place RED probe on (+), BLACK probe on (-)

Expected reading: 16.8V (fully charged 4S)\\
Safe minimum: 13.2V (3.3V per cell $\times$ 4)
\end{examplebox}

\begin{examplebox}
\textbf{EXAMPLE 2: Checking BEC Output}

BEC output connector:\\
Red wire (+5V) - RED probe\\
Black wire (GND) - BLACK probe

Expected reading: $5.0V \pm 0.1V$\\
If reading $< 4.8V \rightarrow$ Problem! BEC may be overloaded or failing
\end{examplebox}

\begin{examplebox}
\textbf{EXAMPLE 3: Voltage Drop Across Wire}

Sometimes you want to check if a wire has too much resistance:

Battery (+) at Point A: 14.8V\\
Component (+) at Point B: 14.5V\\
Voltage drop: 0.3V $\leftarrow$ This is lost as heat in the wire!

$\rightarrow$ If drop $>$ 0.5V, use thicker wire gauge
\end{examplebox}

\begin{warningbox}
SAFETY NOTE: Never measure voltage with power off—it will read 0!
\end{warningbox}

\subsection{Measuring Current}

\begin{warningbox}
MOST DANGEROUS MEASUREMENT - CAN BLOW FUSE OR DAMAGE MULTIMETER!
\end{warningbox}

Current is measured IN SERIES (breaking the circuit)

\textbf{PROCEDURE:}
\begin{enumerate}
\item Turn OFF power first!
\item Move RED probe to ``A'' or ``10A'' jack (check meter's max current!)
\item Set dial to ``A'' (current)
\item BREAK the circuit at one point
\item Connect multimeter IN the current path (multimeter completes the circuit)
\item Turn power ON
\item Current flows THROUGH the meter - read display
\item Turn power OFF before disconnecting!
\end{enumerate}

\begin{warningbox}
NEVER CONNECT MULTIMETER IN PARALLEL IN CURRENT MODE!

This creates a short circuit $\rightarrow$ blown fuse or damaged meter
\end{warningbox}

\begin{examplebox}
\textbf{EXAMPLE: Measuring Servo Current}

You want to know how much current a servo draws:

Original: UBEC (+5V) --- Servo Red Wire

Modified: UBEC (+5V) --- Multimeter RED (A jack) --- Multimeter BLACK --- Servo Red Wire

Ground connection unchanged: UBEC GND --- Servo Black Wire

Turn on power, move servo through full range:
\begin{itemize}[nosep]
\item Reading @ rest: 50 mA
\item Reading @ moving: 300 mA
\item Reading @ stall: 800 mA $\leftarrow$ This is what you need for budgeting!
\end{itemize}
\end{examplebox}

\textbf{SAFER ALTERNATIVE:} Use a ``current sensor module'' or ``power analyzer''
\begin{itemize}
\item Inline device that measures current without breaking circuit
\item Safer for beginners
\item Can log current over time
\end{itemize}

\subsection{Measuring Resistance (and Continuity)}

Resistance is measured with POWER OFF

\textbf{PROCEDURE:}
\begin{enumerate}
\item Disconnect component from circuit (or at least one end)
\item Set dial to ``$\Omega$'' (ohms)
\item Select range (200$\Omega$, 2k$\Omega$, 20k$\Omega$, etc.)
\item Touch probes to component leads
\item Read resistance value
\end{enumerate}

\textbf{WHY DISCONNECT FROM CIRCUIT?}\\
Other components in parallel will affect the reading!

\begin{examplebox}
\textbf{EXAMPLE 1: Testing a Resistor}

Resistor labeled ``220$\Omega$'':
\begin{itemize}[nosep]
\item Measured: 217$\Omega$ $\leftarrow$ Within tolerance (usually $\pm$5\%)
\end{itemize}

Resistor labeled ``1k$\Omega$'' (1000$\Omega$):
\begin{itemize}[nosep]
\item Measured: 1.02k$\Omega$ $\leftarrow$ Good!
\end{itemize}
\end{examplebox}

\begin{examplebox}
\textbf{EXAMPLE 2: Testing Wire Continuity}

Most useful multimeter function for UAV assembly!

Set dial to continuity mode (beeper icon)

Touch probes to both ends of a wire:
\begin{itemize}[nosep]
\item BEEP + Low resistance ($< 1\Omega$) $\rightarrow$ Good connection
\item No beep + ``OL'' (overload) $\rightarrow$ Broken wire!
\end{itemize}

\textbf{PRACTICAL USE: After soldering connectors}
\begin{itemize}[nosep]
\item Touch RED to pin 1 on one connector
\item Touch BLACK to pin 1 on other end
\item Should beep $\rightarrow$ Verified connection
\item Repeat for all pins
\end{itemize}
\end{examplebox}

\begin{examplebox}
\textbf{EXAMPLE 3: Checking for Short Circuits}

BEFORE connecting battery:
\begin{itemize}[nosep]
\item Set multimeter to continuity mode
\item Touch RED to battery (+) connector
\item Touch BLACK to battery (-) connector
\end{itemize}

Should NOT beep! (Resistance should be high)

If it beeps $\rightarrow$ SHORT CIRCUIT! Do not connect battery! Something is connecting (+) and (-) directly
\end{examplebox}

\subsection{Multimeter Troubleshooting Checklist}

\textbf{PROBLEM: ``My flight controller won't turn on''}

\textbf{Step 1:} Measure voltage AT the BEC output
\begin{itemize}[nosep]
\item Expected: $\sim$5.0V
\item If 0V $\rightarrow$ BEC is dead or battery is dead
\item If correct $\rightarrow$ Problem is downstream
\end{itemize}

\textbf{Step 2:} Measure voltage AT the flight controller power pins
\begin{itemize}[nosep]
\item Expected: $\sim$5.0V
\item If 0V but BEC shows 5V $\rightarrow$ Bad wire or connector
\item If correct $\rightarrow$ Flight controller is faulty
\end{itemize}

\textbf{Step 3:} Check current draw (advanced)
\begin{itemize}[nosep]
\item If BEC voltage drops under load $\rightarrow$ Overload or short circuit
\end{itemize}

\vspace{1em}

\textbf{PROBLEM: ``Servo twitches but won't move''}

\textbf{Step 1:} Measure servo power supply voltage WHILE servo is moving
\begin{itemize}[nosep]
\item Expected: $\sim$5.0V
\item If voltage drops to 4.0V $\rightarrow$ Insufficient current, UBEC overloaded
\item If stable at 5.0V $\rightarrow$ Servo may be mechanically jammed
\end{itemize}

\textbf{Step 2:} Measure current draw during movement
\begin{itemize}[nosep]
\item If $>$ UBEC rating $\rightarrow$ Need higher-current UBEC
\end{itemize}

\vspace{1em}

\textbf{PROBLEM: ``GPS won't get satellite lock''}

\textbf{Step 1:} Multimeter can't help much here, but check basics:
\begin{itemize}[nosep]
\item Voltage at GPS power pins: $\sim$5.0V
\item Continuity of TX/RX signal wires from GPS to flight controller
\end{itemize}

Usually not electrical—check:
\begin{itemize}[nosep]
\item GPS has clear view of sky
\item Antenna oriented correctly
\item Wait 5+ minutes on first power-up
\end{itemize}

\newpage
\section{Understanding PWM Signals}

\subsection{What is PWM?}

PWM = Pulse Width Modulation

A digital signal that controls servos, ESCs, and other devices by varying the WIDTH of pulses, not the voltage level.

\textbf{SERVO/ESC PWM SIGNAL CHARACTERISTICS:}
\begin{itemize}
\item Voltage: 5V (or 3.3V for some flight controllers)
\item Frequency: 50 Hz (one pulse every 20 milliseconds)
\item Pulse width: 1000 $\mu$s to 2000 $\mu$s (1 ms to 2 ms)
\end{itemize}

\textbf{VISUALIZING PWM:}

\begin{verbatim}
     ┌───┐                     ┌───┐                     ┌───┐
     │   │                     │   │                     │   │
5V ──┘   └─────────────────────┘   └─────────────────────┘   └────

     ↑                         ↑
     │←─ 1 ms ────→│          (Pulse repeats every 20 ms)
     
     Duty cycle = 1 ms / 20 ms = 5%
\end{verbatim}

\subsection{PWM Control Values}

\textbf{SERVO POSITION:}

\begin{center}
\begin{tabular}{ll}
\textbf{Pulse Width} & \textbf{Servo Position} \\
\hline
1000 $\mu$s & 0$^\circ$ (full left) \\
1500 $\mu$s & 90$^\circ$ (center) \\
2000 $\mu$s & 180$^\circ$ (full right) \\
\end{tabular}
\end{center}

\textbf{THROTTLE/ESC:}

\begin{center}
\begin{tabular}{ll}
\textbf{Pulse Width} & \textbf{Throttle} \\
\hline
1000 $\mu$s & 0\% (motor off) \\
1500 $\mu$s & 50\% throttle \\
2000 $\mu$s & 100\% throttle (full power) \\
\end{tabular}
\end{center}

\begin{warningbox}
NOTE: Some ESCs use different ranges (e.g., 1100-1900 $\mu$s). Check your ESC datasheet and calibrate!
\end{warningbox}

\subsection{How to Measure PWM with Multimeter}

\begin{warningbox}
Regular multimeters CANNOT accurately measure PWM pulse width!
\end{warningbox}

However, you CAN verify signal presence:

\begin{enumerate}
\item Set multimeter to DC voltage
\item Touch RED to signal wire, BLACK to ground
\item Reading will show AVERAGE voltage

Example:
\begin{itemize}[nosep]
\item 1 ms pulse / 20 ms period = 5\% duty cycle
\item Average voltage = 5V $\times$ 0.05 = 0.25V
\item 2 ms pulse / 20 ms period = 10\% duty cycle
\item Average voltage = 5V $\times$ 0.10 = 0.5V
\end{itemize}

\item If you see $\sim$0.25V to 0.5V $\rightarrow$ PWM signal is present
\item If you see 0V or 5V constantly $\rightarrow$ Problem!
\end{enumerate}

\textbf{TO ACCURATELY MEASURE PWM:}
\begin{itemize}
\item Use an oscilloscope (expensive!)
\item Use a servo tester (cheap device that displays pulse width)
\item Use Arduino sketch to read pulseIn() function
\end{itemize}

\subsection{PWM Troubleshooting}

\textbf{PROBLEM: ``ESC won't arm / motor won't spin''}

ESCs require a specific startup sequence:
\begin{enumerate}[nosep]
\item Throttle stick at minimum (1000 $\mu$s)
\item Connect battery
\item Wait for ESC beeps (indicates cells detected)
\item Raise throttle slightly $\rightarrow$ Motor should spin
\end{enumerate}

If motor doesn't respond:
\begin{itemize}[nosep]
\item Check PWM signal wire is connected (usually white/yellow)
\item Verify ESC is receiving power (check voltage)
\item Calibrate ESC throttle range (see ESC manual)
\end{itemize}

\vspace{1em}

\textbf{PROBLEM: ``Servo jitters or moves erratically''}

\textbf{CAUSES:}
\begin{enumerate}
\item \textbf{Electrical noise on signal wire}
\begin{itemize}[nosep]
\item Solution: Twist signal wire with ground wire
\item Keep signal wires away from motor wires
\end{itemize}

\item \textbf{Insufficient power (voltage drop)}
\begin{itemize}[nosep]
\item Measure voltage at servo while it moves
\item If drops below 4.5V $\rightarrow$ Need thicker power wires or higher-current BEC
\end{itemize}

\item \textbf{Bad connection}
\begin{itemize}[nosep]
\item Check continuity of all three servo wires
\item Re-solder connectors
\end{itemize}
\end{enumerate}

\newpage
\section{Grounding and Electrical Noise}

\subsection{Why Grounding Matters}

GROUND = REFERENCE POINT for all voltages in your system

Without a proper common ground:
\begin{itemize}
\item Voltages are undefined
\item Signals won't be read correctly
\item Electrical noise causes interference
\item Components may be damaged
\end{itemize}

\begin{keyinsight}
\textbf{GOLDEN RULE:}\\
ALL COMPONENTS IN YOUR UAV MUST SHARE A COMMON GROUND!
\end{keyinsight}

\subsection{Ground Loops and Noise}

\textbf{GROUND LOOP:} Multiple paths between grounds create circulating currents

\textbf{EXAMPLE OF PROBLEM:}

If Ground wire A has different resistance than Ground wire B, a small voltage difference exists between the two ``grounds''!

This can cause:
\begin{itemize}
\item ADC reading errors on Arduino
\item PWM signal corruption
\item Random resets or glitches
\end{itemize}

\textbf{SOLUTION: Star grounding}

All grounds connect to ONE central point (usually battery negative):

\begin{center}
Battery (-) $\leftarrow$ Central ground point\\
$\downarrow$\\
All components connect here
\end{center}

\subsection{Sources of Electrical Noise}

\textbf{1. MOTOR/ESC NOISE}

Brushless motors create high-frequency electrical noise

Effects:
\begin{itemize}
\item Corrupts GPS reception
\item Causes servo jitter
\item Interferes with radio receiver
\end{itemize}

\textbf{SOLUTIONS:}
\begin{itemize}
\item Keep motor wires short
\item Twist motor wires together (reduces electromagnetic radiation)
\item Add capacitors across motor terminals (suppresses noise)
\item Separate signal wires from motor wires (at least 5 cm apart)
\end{itemize}

\textbf{2. POWER SUPPLY NOISE}

When servos move, they draw sudden current spikes\\
This causes voltage ripple on power rails

\textbf{SOLUTIONS:}
\begin{itemize}
\item Add capacitors near each component (decoupling capacitors)
\begin{itemize}
\item 100 $\mu$F electrolytic capacitor on BEC output
\item 0.1 $\mu$F ceramic capacitor at each component's power pins
\end{itemize}
\item Use separate BECs for servos vs. flight-critical systems (ALREADY DOING THIS!)
\end{itemize}

\textbf{3. RADIO FREQUENCY INTERFERENCE (RFI)}

Long wires act as antennas, picking up noise

\textbf{SOLUTIONS:}
\begin{itemize}
\item Keep wires as short as practical
\item Twist signal + ground pairs together
\item Use shielded cable for sensitive signals (e.g., GPS antenna)
\end{itemize}

\subsection{Best Practices for Your UAV}

\begin{itemize}
\item CONNECT all component grounds to battery negative terminal (directly or through BECs that connect to battery)
\item TWIST signal wires with their ground/return wires
\begin{itemize}
\item Reduces crosstalk between signals
\item Minimizes antenna effect (noise pickup)
\end{itemize}
\item ROUTE motor wires AWAY from signal wires
\begin{itemize}
\item If they must cross, do so at 90$^\circ$ angle (minimizes coupling)
\end{itemize}
\item ADD a large capacitor (100-220 $\mu$F) across battery terminals
\begin{itemize}
\item Smooths voltage spikes when motor draws high current
\end{itemize}
\item TEST grounding by measuring voltage between grounds
\begin{itemize}
\item Connect multimeter BLACK to battery (-)
\item Touch RED to ground pin of each component
\item Should read 0.00V (or $< 0.01V$)
\item If $> 0.1V \rightarrow$ Poor ground connection!
\end{itemize}
\end{itemize}

\newpage
\section{Systematic Troubleshooting Approach}

\subsection{The Troubleshooting Mindset}

When something doesn't work, DON'T just try random fixes!

Follow this systematic approach:
\begin{enumerate}
\item \textbf{OBSERVE} - What exactly is wrong? Be specific.
\item \textbf{ISOLATE} - Which component or connection is faulty?
\item \textbf{VERIFY} - Use measurements (multimeter) to confirm your hypothesis
\item \textbf{FIX} - Make ONE change at a time
\item \textbf{TEST} - Did the fix work?
\end{enumerate}

\subsection{Power Troubleshooting Flowchart}

\textbf{SYMPTOM: ``Nothing turns on''}

\begin{enumerate}
\item \textbf{Check battery voltage (with multimeter)}
\begin{itemize}
\item Is it $> 13.2V$? (3.3V/cell minimum)
\item NO $\rightarrow$ Charge battery
\item YES $\rightarrow$ Continue
\end{itemize}

\item \textbf{Check XT60 connector continuity}
\begin{itemize}
\item Are positive and negative both connected?
\item NO $\rightarrow$ Resolder connector
\item YES $\rightarrow$ Continue
\end{itemize}

\item \textbf{Measure voltage at ESC input (14.8V expected)}
\begin{itemize}
\item 0V $\rightarrow$ Bad Y-harness or connector
\item 14.8V $\rightarrow$ Continue
\end{itemize}

\item \textbf{Measure voltage at ESC BEC output (5V expected)}
\begin{itemize}
\item 0V $\rightarrow$ ESC is dead or BEC damaged, try external UBEC alone
\item 5V $\rightarrow$ Continue
\end{itemize}

\item \textbf{Measure voltage at component power pins}
\begin{itemize}
\item 0V $\rightarrow$ Bad wire from BEC to component, resolder
\item 5V $\rightarrow$ Component is faulty
\end{itemize}
\end{enumerate}

\subsection{Component-Specific Troubleshooting}

\textbf{FLIGHT CONTROLLER WON'T BOOT:}
\begin{enumerate}
\item Verify 5V on power pins (multimeter)
\item Check current draw (should be $\sim$100-200 mA)
\begin{itemize}[nosep]
\item If 0 mA $\rightarrow$ No power reaching flight controller
\item If $> 500$ mA $\rightarrow$ Possible short circuit
\end{itemize}
\item Check USB connection (can it connect to computer?)
\begin{itemize}[nosep]
\item YES $\rightarrow$ Power issue, check external power supply
\item NO $\rightarrow$ Flight controller may be bricked, reflash firmware
\end{itemize}
\end{enumerate}

\textbf{GPS NO SATELLITE LOCK:}
\begin{enumerate}
\item Verify 5V on GPS power pins
\item Check continuity of TX/RX wires to flight controller
\item Verify GPS is configured for correct baud rate (usually 38400)
\item Test outdoors with clear sky view (takes 5-15 min on first boot)
\item Check for electrical noise:
\begin{itemize}[nosep]
\item Move GPS farther from motor/ESC
\item Add aluminum foil shield under GPS (acts as ground plane)
\end{itemize}
\end{enumerate}

\textbf{SERVO NOT RESPONDING:}
\begin{enumerate}
\item Verify 5V on servo power pins while it attempts to move
\begin{itemize}[nosep]
\item If drops $< 4.5V \rightarrow$ Insufficient current (upgrade UBEC)
\end{itemize}
\item Test servo with servo tester (bypass flight controller)
\begin{itemize}[nosep]
\item Works $\rightarrow$ Problem is flight controller or wiring
\item Doesn't work $\rightarrow$ Servo is dead
\end{itemize}
\item Check PWM signal continuity from flight controller to servo
\item Verify flight controller channel is configured correctly (transmitter mixer)
\end{enumerate}

\textbf{MOTOR WON'T SPIN:}
\begin{enumerate}
\item Verify 14.8V at ESC input
\item Check ESC programming (throttle calibration)
\item Test motor directly with servo tester (if 3-wire ESC)
\item Verify motor isn't mechanically jammed (try spinning by hand)
\item Check for loose motor wires (intermittent connection)
\end{enumerate}

\textbf{ARDUINO NOT LOGGING DATA:}
\begin{enumerate}
\item Verify 5V power at Arduino
\item Test sensors individually (I2C scanner sketch)
\item Check SD card:
\begin{itemize}[nosep]
\item Is it formatted FAT32?
\item Is it inserted correctly?
\item Test write with simple Arduino sketch
\end{itemize}
\item Add Serial.print() debugging to code to see where it fails
\end{enumerate}

\subsection{Common Mistakes to Avoid}

\begin{enumerate}
\item \textbf{MIXING UP VOLTAGE LEVELS}
\begin{itemize}[nosep]
\item Connecting 5V to 3.3V component $\rightarrow$ Damage!
\item Check datasheet for EVERY component
\end{itemize}

\item \textbf{REVERSING POLARITY}
\begin{itemize}[nosep]
\item Connecting battery backwards $\rightarrow$ EXPLOSION (LiPo fire!)
\item Double-check red = (+), black = (-) BEFORE connecting power
\end{itemize}

\item \textbf{FORGETTING COMMON GROUND}
\begin{itemize}[nosep]
\item Connecting only (+) and signal, forgetting ground wire
\item Every component needs THREE connections: +, signal, ground
\end{itemize}

\item \textbf{OVERLOADING BEC}
\begin{itemize}[nosep]
\item Connecting too many servos to one BEC
\item Calculate total current draw FIRST, compare to BEC rating
\end{itemize}

\item \textbf{MEASURING CURRENT IN PARALLEL}
\begin{itemize}[nosep]
\item Multimeter in current mode connected like a voltmeter $\rightarrow$ BOOM!
\item Current is measured IN SERIES (break the circuit)
\end{itemize}

\item \textbf{TESTING WITH PROPELLER ATTACHED}
\begin{itemize}[nosep]
\item Motor spins unexpectedly $\rightarrow$ Injury!
\item ALWAYS remove propeller during electrical testing
\end{itemize}
\end{enumerate}

\newpage
\section{Practical Exercises}

\subsection*{EXERCISE 1: Datasheet Scavenger Hunt}

Find datasheets for three components in your UAV system:
\begin{enumerate}
\item Your ESC
\item Your Arduino/microcontroller
\item One sensor (BME280, air quality sensor, etc.)
\end{enumerate}

For each datasheet, extract and record:
\begin{itemize}
\item Operating voltage range
\item Current consumption (typical and maximum)
\item Pinout diagram (draw it or save screenshot)
\item I2C address (if applicable)
\item One ``interesting fact'' from the datasheet you didn't know
\end{itemize}

\textbf{DUE DATE:} Before next meeting

\subsection*{EXERCISE 2: Multimeter Practice}

Perform these measurements and record results:

\textbf{1. Battery voltage:}\\
Measured: \underline{\hspace{3cm}} V\\
Calculated per-cell voltage: \underline{\hspace{3cm}} V (divide by 4)\\
Status: $\square$ Good ($>3.7V$/cell) \quad $\square$ Needs charge \quad $\square$ Dead

\textbf{2. USB phone charger output:}\\
Measured: \underline{\hspace{3cm}} V (should be $\sim$5V)

\textbf{3. Resistance of 1-meter wire:}\\
Wire gauge: \underline{\hspace{3cm}} AWG\\
Measured: \underline{\hspace{3cm}} $\Omega$ (should be $< 1\Omega$ for short wires)

\textbf{4. Continuity test on a homemade cable:}\\
Wire 1: $\square$ Beeps (good) \quad $\square$ No beep (broken)\\
Wire 2: $\square$ Beeps (good) \quad $\square$ No beep (broken)\\
Wire 3: $\square$ Beeps (good) \quad $\square$ No beep (broken)

\textbf{DUE DATE:} Complete during next lab session

\subsection*{EXERCISE 3: Draw Your System Schematic}

Create a schematic diagram of YOUR UAV electrical system showing:
\begin{itemize}
\item Battery (with voltage label)
\item Y-harness
\item ESC and external UBEC
\item Both power rails (Rail 1 and Rail 2)
\item All components: Flight controller, Arduino, GPS, RC receiver, servos
\item Ground connections (common ground symbol)
\end{itemize}

You can draw by hand or use:
\begin{itemize}
\item Draw.io (free, online)
\item Fritzing (beginner-friendly)
\item KiCad (professional, but complex)
\end{itemize}

\textbf{HINT:} Reference the schematic example in Section 1.4

\textbf{DUE DATE:} Before flight testing begins

\subsection*{EXERCISE 4: PWM Investigation}

Write a simple Arduino sketch that:
\begin{enumerate}
\item Reads PWM input from RC receiver on pin 2
\item Prints the pulse width (in microseconds) to Serial Monitor
\item Maps the pulse width to a percentage (1000 = 0\%, 2000 = 100\%)
\end{enumerate}

Use this to verify your RC receiver outputs correct PWM values.

\textbf{TEST:}\\
Throttle stick at minimum $\rightarrow$ Should read $\sim$1000 $\mu$s\\
Throttle stick at center $\rightarrow$ Should read $\sim$1500 $\mu$s\\
Throttle stick at maximum $\rightarrow$ Should read $\sim$2000 $\mu$s

\textbf{Record actual measured values:}\\
Minimum: \underline{\hspace{3cm}} $\mu$s\\
Center: \underline{\hspace{3cm}} $\mu$s\\
Maximum: \underline{\hspace{3cm}} $\mu$s

\textbf{ARDUINO CODE HINT:}
\begin{verbatim}
int pulse = pulseIn(2, HIGH, 25000);  // Read pulse on pin 2
Serial.println(pulse);
\end{verbatim}

\textbf{DUE DATE:} Whenever you receive your RC equipment

\newpage
\section{Advanced Topic - Decoupling Capacitors}

\textit{(This section is OPTIONAL but highly recommended for understanding why some circuits include capacitors you didn't think were necessary!)}

\subsection{Why Decoupling Capacitors?}

When a component (like a servo or Arduino) suddenly draws current, the voltage on the power rail can temporarily drop. This is called a ``voltage sag'' or ``brownout.''

\textbf{CAUSES OF VOLTAGE SAG:}
\begin{enumerate}
\item Wire resistance (even tiny amounts matter!)
\item BEC can't respond instantly to current changes
\item Inductance in long wires (resist changes in current)
\end{enumerate}

\textbf{SOLUTION:} Capacitor acts as a ``local battery''
\begin{itemize}
\item Stores charge near the component
\item Supplies current during sudden demand
\item Smooths voltage ripple
\end{itemize}

\textbf{ANALOGY: Water tower near your house}
\begin{itemize}
\item City water supply = BEC
\item Pipes = Wires (have resistance)
\item Water tower = Capacitor (local storage)
\end{itemize}

When you suddenly turn on all faucets, the water tower supplies water instantly while the city pumps catch up.

\subsection{Types of Capacitors}

\textbf{ELECTROLYTIC CAPACITOR (100 $\mu$F - 1000 $\mu$F)}
\begin{itemize}
\item Large capacitance
\item Good for bulk energy storage
\item Polarized (must connect (+) to (+), (-) to (-))
\item Slower response time
\end{itemize}

\textbf{USE:} At BEC outputs, near power-hungry components

\textbf{CERAMIC CAPACITOR (0.1 $\mu$F - 10 $\mu$F)}
\begin{itemize}
\item Small capacitance
\item Very fast response time
\item Not polarized (can connect either way)
\item Physically small
\end{itemize}

\textbf{USE:} Right next to IC chips (flight controller, Arduino)

\subsection{Where to Add Capacitors in Your UAV}

\textbf{1. ACROSS BATTERY TERMINALS}
\begin{itemize}
\item 220 $\mu$F to 470 $\mu$F electrolytic
\item Reduces voltage spikes when motor draws high current
\item Protects ESC and BECs
\end{itemize}

\textbf{2. AT EACH BEC OUTPUT}
\begin{itemize}
\item 100 $\mu$F electrolytic + 0.1 $\mu$F ceramic
\item Smooths BEC output ripple
\item Place physically close to BEC
\end{itemize}

\textbf{3. AT EACH COMPONENT'S POWER PINS}
\begin{itemize}
\item 0.1 $\mu$F ceramic
\item As close as possible to IC power pins
\item Flight controller and Arduino especially benefit
\end{itemize}

\textbf{CONNECTION:}\\
Always connect capacitor between (+) and (-) [or (+) and GND]

\newpage
\section{Preparing for Hands-On Assembly}

\subsection{Pre-Assembly Checklist}

BEFORE you start assembling:

\begin{itemize}[label=$\square$]
\item All components received and inspected
\item Datasheets downloaded and key specs recorded in Quick Reference Card
\item Multimeter tested and working (check battery, test on known voltage)
\item Soldering iron heated and tip tinned
\item Wiring diagram printed and accessible
\item Safety Protocol reviewed (especially LiPo handling!)
\item Work area clean and organized
\item Fire extinguisher or water nearby (LiPo safety)
\end{itemize}

\subsection{Assembly Sequence (Recommended)}

\textbf{PHASE 1: Power System (DO THIS FIRST)}
\begin{enumerate}
\item Install XT60 connectors on battery and Y-harness
\item Test continuity of Y-harness
\item Connect ESC and UBEC to Y-harness
\item Power on (no other components connected yet)
\item Verify ESC BEC outputs 5V
\item Verify UBEC outputs 5V
\item Measure current draw (should be very low, $\sim$10-50 mA)
\end{enumerate}

\begin{warningbox}
DO NOT PROCEED if voltages are incorrect!
\end{warningbox}

\textbf{PHASE 2: Flight-Critical Components (RAIL 1)}
\begin{enumerate}
\item Connect flight controller to ESC BEC
\item Verify flight controller boots (LEDs light up)
\item Connect RC receiver
\item Verify RC receiver binds to transmitter
\item Connect GPS module
\item Wait for GPS satellite lock
\end{enumerate}

\textbf{Test:} All components powered simultaneously, battery voltage $> 14V$

\textbf{PHASE 3: Data Collection System}
\begin{enumerate}
\item Connect Arduino to RAIL 1
\item Upload test sketch (blink LED or print to serial)
\item Connect sensors one at a time
\item Test each sensor individually
\item Install SD card module and test data logging
\end{enumerate}

\textbf{Test:} Log data for 5 minutes, verify file saved correctly

\textbf{PHASE 4: Servos (RAIL 2)}
\begin{enumerate}
\item Connect ONE servo to external UBEC
\item Verify servo moves through full range
\item Measure current draw at stall
\item Connect remaining servos
\item Test all servos simultaneously
\item Verify UBEC voltage doesn't drop below 4.8V under load
\end{enumerate}

\textbf{Test:} All servos to full deflection at once, no brownouts

\textbf{PHASE 5: Motor System (SAVE FOR LAST)}
\begin{enumerate}
\item Connect motor to ESC (NO PROP YET!)
\item Connect ESC signal wire to flight controller
\item Arm ESC (throttle low, then battery, wait for beeps)
\item Test motor at 10\% throttle (should spin smoothly)
\item Measure motor current
\item If all good, install propeller (secured drone first!)
\end{enumerate}

\subsection{Testing Before Flight}

\textbf{BENCH TEST (no propeller, secured):}
\begin{itemize}[label=$\square$]
\item All components powered on
\item Flight controller boots and shows GPS lock
\item RC transmitter controls all servos correctly
\item Throttle stick moves (but ESC disarmed for safety)
\item Data logging works continuously for 10 minutes
\item Battery voltage after 10 min test: \underline{\hspace{2cm}} V
\end{itemize}

\textbf{GROUND TEST (propeller on, drone secured):}
\begin{itemize}[label=$\square$]
\item Arm ESC
\item Slowly increase throttle to 50\%
\item Verify motor spins in correct direction
\item Check for vibrations (balance propeller if needed)
\item Measure current draw: \underline{\hspace{2cm}} A
\item Verify control surfaces respond correctly
\item Test failsafe (turn off transmitter $\rightarrow$ motor stops)
\end{itemize}

\textbf{RANGE TEST:}
\begin{itemize}[label=$\square$]
\item Walk 50 meters away with transmitter
\item Verify all controls still respond
\item Repeat at 100 meters
\item If control is lost, do NOT fly! Check antenna placement.
\end{itemize}

\newpage
\section{Summary and Next Steps}

\subsection*{What You Learned in This Lesson:}

\begin{itemize}
\item How to read electrical schematics and wiring diagrams
\item How to extract critical information from component datasheets
\item How to use a multimeter to measure voltage, current, and continuity
\item How PWM signals work and how they control servos/ESCs
\item Why grounding is critical and how to avoid ground loops
\item Systematic troubleshooting approach for electrical problems
\item Where to add capacitors to reduce noise and voltage sags
\end{itemize}

\subsection*{Skills You Can Now Apply:}

\begin{itemize}
\item Debug ``component won't turn on'' problems using multimeter
\item Verify correct voltage at each point in your system
\item Test continuity of all connections before powering on
\item Identify and fix grounding issues
\item Read your component datasheets to find pinouts and specs
\item Systematically troubleshoot problems instead of guessing
\end{itemize}

\subsection*{What's Next? (Lesson 3 Preview)}

The next lesson will cover:

\textbf{LESSON 3: SOLDERING, CONNECTORS, AND PROFESSIONAL ASSEMBLY TECHNIQUES}

Topics:
\begin{itemize}
\item How to solder properly (technique, temperature, inspection)
\item Choosing the right connectors (XT60, JST, Dupont, bullet connectors)
\item Crimping vs. soldering
\item Heat shrink tubing and strain relief
\item Wire management and cable routing
\item Building a custom wiring harness
\item Quality control and inspection
\item Common soldering mistakes and how to fix them
\end{itemize}

\subsection*{Homework Before Next Lesson:}

\begin{enumerate}
\item Complete Exercise 1: Datasheet Scavenger Hunt
\item Fill out Component Quick Reference Card (as components arrive)
\item Practice with multimeter (Exercise 2)
\item Review soldering safety videos online (search ``soldering basics'')
\end{enumerate}

\subsection*{Questions for Review:}

Test your understanding:

\begin{enumerate}
\item You measure 5.2V at the BEC output, but only 4.5V at the flight controller. What is the most likely cause, and how do you fix it?

\item Your servo twitches randomly. List three possible causes and how you would test for each one.

\item Why is it dangerous to measure current with the multimeter set incorrectly?

\item Your motor spins, but GPS won't get satellite lock. Could these be related? Explain.

\item What is the purpose of connecting all grounds to a common point?
\end{enumerate}

(Answers will be discussed in next lesson)

\vspace{2em}

\begin{center}
\rule{\textwidth}{2pt}
\textbf{\Large CONGRATULATIONS!}\\[0.5em]
\rule{\textwidth}{2pt}
\end{center}

You now have the foundational knowledge to read schematics, interpret datasheets, and troubleshoot electrical problems systematically.

These skills will serve you throughout this project and in future electrical engineering work.

\begin{center}
\textbf{\Large REMEMBER: When in doubt, measure first, then act!}
\end{center}

\vspace{2em}

\begin{center}
\rule{\textwidth}{2pt}
\textbf{END OF LESSON 2}\\
\rule{\textwidth}{2pt}
\end{center}

\end{document}
